% Based on https://github.com/JavierGelatti/Formato-LaTeX-CoNaIISI

\input{format.tex}

\author{Castaño Ezequiel Leonardo}
\title{Presentación de Trabajos Estudiantiles}
\email{skielcast@gmail.com}
\universidad{Universidad Tecnológica Nacional Facultad Regional Rosario}

\begin{document}
\titulo

\begin{multicols}{2}

\resumen{
Un abstract tiene típicamente un único párrafo y menos de 250 palabras y debe permitir a los lectores identificar el contenido básico del documento rápida y fielmente, con el fin de determinar la relevancia del mismo para sus intereses y, por tanto, para decidir si necesitan leer el documento en su totalidad. No debe tener citas bibliográficas. Debería responder, entre otras,  a las siguientes preguntas: ¿Por qué se hizo?, ¿Qué hizo?, ¿Cómo  lo hizo?, ¿Cuáles serán los beneficios? (Times New Roman, 10, cursiva) 
}

\miniseccion{Palabras Clave}

{\footnotesize
Ayudan a localizar en forma rápida y eficiente el artículo correcto en bases de datos científicas. Son los términos más relevantes que tienen desarrollo en el texto. (Times New Roman, 10).
\par}

\section{Introducción}

La introducción debe responder a la pregunta: “¿por qué se ha hecho este trabajo?". Describe el interés que tiene en el contexto científico/académico del momento. Se deben demostrar, a través del análisis crítico, que existe un problema que amerita ser abordado. Debería contestar (entre otras) a las preguntas: ¿Qué se conoce sobre el desarrollo o el tema investigado? ¿Cuál es el problema que se quiere resolver? ¿Cuál es la utilidad que se espera de la solución? Y ¿Qué aplicaciones prácticas tendrá? Describir la organización del trabajo. (Times New Roman, 12).

\section{Elementos del Trabajo y Metodología}

Es la parte sustancial del trabajo, el lector debe comprender el método usado con tal detalle que le permita aplicarlo al mismo o a otro problema. Se desarrollan los conceptos explicando la metodología utilizada. Será escrito en fuente.

\section{Resultados}

Los resultados son los que avalarán las conclusiones y justificarán la utilidad el trabajo realizado.

\section{Trabajos Relacionados}

En esta sección se deben detallar, con sentido crítico, trabajos vinculados realizados y cómo se ha resuelto este problema (o problemas similares) en otros trabajos, si los hubiera. (Times New Roman, 12).

\section{Conclusión y Trabajos Futuros}

Esta sección debe explicitar las limitaciones del trabajo presentado y establecer una discusión sobre los resultados o conclusiones presentadas. Se debe realizar un análisis de los aportes del trabajo frente a otros anteriores si los hubiera. Además, se deben establecer cuestiones abiertas y probables líneas adicionales en el marco de los resultados obtenidos. Los trabajos futuros se deben relacionar con la superación de las limitaciones del trabajo presentado. Muchas veces es, junto con el título, la parte más leída y por lo tanto debe ser de fácil comprensión. (Times New Roman, 12).

\section{Referencias}

Documentación y bibliografía utilizada. Todas las publicaciones citadas deberán incluirse en la lista de referencias. La numeración será secuencial y estará entre corchetes: [1]. Será escrita en fuente (Times New Roman, 10)

\input{bibliography.tex}

\miniseccion{Datos de Contacto}

{\it\footnotesize

Nombre y Apellido. Institución. Dirección postal. Email. 

\par}


Tablas: Las Tablas serán numeradas en forma secuencial (Tabla 1, 2, 3, etc.,) con un titulo descriptivo, ambos escritos en Times New Roman, 10, cursiva. Las tablas estarán centradas en la columna o en caso contrario en toda la página. Los títulos de las columnas de la Tabla estarían en caracteres Times New Roman, 10 negrita.


Figuras: Las figuras serán numeradas secuencialmente: Figura 1, 2, 3, etc.).
Aquellas figuras que abarquen toda la página serán ubicadas al comienzo o al final de la página.


Es preferible que los dibujos, fotos, diagramas, etc., se incluyan en blanco y negro. De esta manera, las impresiones monocromáticas, representarán fielmente el documento original.

\end{multicols}

\end{document}